\section{Project Timeline and Gantt Chart}

Figure~\ref{fig:gantt} shows the schedule the project followed, from data preparation and
exploratory analysis through model training and the statistical study to building the API
and the web application.

\begin{figure}[htbp]
    \centering
    \includegraphics[width=\textwidth]{gantt.png}
    \caption{Project Gantt chart.}
    \label{fig:gantt}
\end{figure}

From the figure above we can see how the phases overlap: modelling work continues while the
API and web application are built, and the statistical study runs after training is
complete.

\section{Supplementary Tables}
\label{app:tables}

\subsection{Full Significance Test Results}

Table~\ref{tab:sig} lists the ten smallest $p$-values. The complete set of 28 pairwise
McNemar comparisons is available in \texttt{reports/\allowbreak significance\_\allowbreak tests.csv}. No comparison
reaches significance at $\alpha = 0.05$ after Holm correction, and the smallest adjusted
value is 0.096.

\subsection{Cost at Every Threshold}

Table~\ref{tab:cost} shows representative operating points. The full 201-row sweep, at
threshold increments of 0.005, is in \texttt{reports/\allowbreak cost\_\allowbreak analysis.csv}.

\subsection{Calibration Bins}

Table~\ref{tab:calibration} reports each model's summary ECE and Brier score. The per-bin
detail (bin boundaries, counts, mean predicted probability, and observed fraud rate) is in
\texttt{reports/\allowbreak calibration.csv}.

\section{Reproducibility}
\label{app:repro}

\subsection{Environment}

All results were produced with Python 3.13 and the following pinned packages:

\begin{center}
\begin{tabular}{@{}ll@{}}
\hline
\textbf{Package} & \textbf{Version} \\
\hline
pandas             & 2.3.3 \\
numpy              & 2.3.5 \\
scikit-learn       & 1.8.0 \\
xgboost            & 3.3.0 \\
lightgbm           & 4.6.0 \\
imbalanced-learn   & 0.14.2 \\
scipy              & 1.17.0 \\
matplotlib         & 3.10.8 \\
joblib             & 1.5.3 \\
fastapi            & 0.128.0 \\
uvicorn            & 0.40.0 \\
\hline
\end{tabular}
\end{center}

Every random component uses \texttt{random\_\allowbreak state = 42}: the train/validation/test splitter,
all resamplers, all classifiers, the bootstrap resampler, the cross-validation fold
assignment, and the permutation-importance shuffles.

\subsection{Reproducing the Results}

% Verbatim never wraps, so every line here is kept under 52 characters and the block is
% set one size down to stay inside the text width.
{\footnotesize
\begin{verbatim}
pip install -r requirements.txt

python src/eda.py            # figures 01-07
python src/train.py          # 8 pipelines, registry,
                             #   figure 08, comparison
python src/resampling_experiment.py
                             # figures 09-10
python src/deep_analysis.py  # figures 11-19, 10 CSVs
python src/make_tables.py    # regenerate all tables

python -m uvicorn api.main:app --port 8000
cd web && npm install && npm run dev
\end{verbatim}
}

\texttt{deep\_\allowbreak analysis.py} accepts \texttt{-{}-fast} to skip cross-validation and the
learning curve, which account for most of its runtime (roughly 7 minutes of the 8-minute
total on the development machine).

\subsection{Verification Built into the Pipeline}

Two automatic checks guard against the errors that would most easily invalidate these
results.

\texttt{train.py} reloads the saved champion from disk, re-scores the test set, and asserts
that the PR-AUC matches the in-memory value to within $10^{-9}$. It also warns if the
PR-AUC exceeds 0.99, which on this dataset would indicate leakage rather than success.

\texttt{deep\_\allowbreak analysis.py} independently reloads all eight pipelines and asserts that the
champion reproduces a PR-AUC of $0.8277$ and confusion counts of $70/2/25/56{,}649$ exactly
before performing any analysis. Every number in Chapter~\ref{ch:results} is therefore
computed from models verified to be identical to those reported in
Section~\ref{sec:comparison}.

\section{API Reference}
\label{app:api}

\subsection{Scoring a Transaction}

{\footnotesize
\begin{verbatim}
POST /predict?model=LightGBM + weighted
             &threshold=0.636
{
  "features": {
    "V1": -1.36, "V2": -0.07, ...,
    "Amount": 149.62
  }
}

# Response
{
  "fraud_probability": 0.0121,
  "is_fraud": false,
  "threshold": 0.636,
  "model": "LightGBM + weighted"
}
\end{verbatim}
}

\subsection{Pipeline Construction}

Each model is one pipeline of scaling, optional resampling, and a classifier. The resampler
is a pipeline step, so it runs only on the training fold during \texttt{fit}. This is the
leakage control described in Section~\ref{sec:leakage}.

{\footnotesize
\begin{verbatim}
for name, clf in classifiers.items():
    models[f"{name} + weighted"] = Pipeline([
        ("scale", StandardScaler()),
        ("clf", clf),
    ])

for name, clf in plain_classifiers.items():
    models[f"{name} + SMOTE"] = Pipeline([
        ("scale", StandardScaler()),
        ("smote", SMOTE(random_state=SEED)),
        ("clf", clf),
    ])
\end{verbatim}
}

\subsection{Threshold Selection}

This is the implementation of Equation~\ref{eq:threshold}, computed in $O(n \log n)$ by
sorting the validation probabilities once and sweeping both error counts with cumulative
sums.

{\footnotesize
\begin{verbatim}
def best_threshold(y_val, val_proba) -> float:
    """Cut-off minimising FN + FP, in O(n log n)."""
    order = np.argsort(val_proba)
    y_sorted = np.asarray(y_val)[order]
    n_pos = int(y_sorted.sum())
    n = len(y_val)

    missed = np.cumsum(y_sorted)
    seen = np.arange(1, n + 1)
    false_alarms = (n - seen) - (n_pos - missed)
    best = int(np.argmin(missed + false_alarms))
    return float(np.nextafter(
        val_proba[order][best], np.inf))
\end{verbatim}
}

\subsection{McNemar's Exact Test}

The code below implements the test of \citet{mcnemar1947note}, in the paired-classifier
form recommended by \citet{dietterich1998approximate}. Only the two discordant counts enter
the test: $b$, the cases model A gets right and model B gets wrong, and $c$, the reverse.
Under the null hypothesis that the two models are equally accurate, $b$ follows a binomial
distribution with $n = b + c$ and $p = 0.5$.

{\footnotesize
\begin{verbatim}
b = int((correct_a & ~correct_b).sum())  # A right
c = int((~correct_a & correct_b).sum())  # B right
p = binomtest(b, b + c, 0.5).pvalue      # exact test
\end{verbatim}
}

The exact binomial form is used rather than the $\chi^2$ approximation because the number of
discordant pairs is between 12 and 17 for the comparisons of interest. That is well below
the count of roughly 25 at which the approximation becomes reliable.
